\section{Controlled Micro-Benchmark Evaluation (RQ2)}
\label{sec:benchmark}

To verify implementation correctness and ensure AST visitors and taint patterns faithfully enforce defined predicates without regression (RQ2), we constructed and open-sourced a suite of 30 compilable micro-testcase applications ($\text{MB-001}$ to $\text{MB-030}$), following established benchmark methodology in Android program analysis~\cite{droidbench_suite, flowdroid}.

\subsection{Methodological Role of Micro-Benchmarks}
\label{sec:benchmark:rationale}

Evaluating static analyzers exclusively on commercial applications introduces confounding factors: obfuscation, packing, and complex dependency graphs obscure whether extraction failures stem from statutory mappings or unrelated decompilation defects. Synthetic micro-testcases provide controlled isolation: each application exercises exactly one security pattern (e.g., \texttt{cleartextTrafficPermitted=false} in $\text{MB-001}$ vs. unencrypted HTTP in $\text{MB-002}$). This enables rigorous verification of positive and negative exemplars for each formalized property.

\begin{table*}[t]
\caption{Implementation Conformance Metrics Across 30 Micro-Applications (RQ2). \textit{Note:} Reports unit-test verification of AST visitors against synthetic exemplars. Wild-world production capability is evaluated in §\ref{sec:ecosystem:adjudication}.}
\label{tab:full_benchmark}
\centering
\small
\resizebox{\textwidth}{!}{%
\begin{tabular}{llccccccc}
\toprule
\textbf{Technical Dimension} & \textbf{Micro-App IDs} & \textbf{Cases ($N$)} & \textbf{Total Cells} & \textbf{Precision} & \textbf{Recall} & \textbf{F1-Score} & \textbf{FPR} & \textbf{FNR} \\
\midrule
\textbf{Network Security (TLS \& Pinning)} & $\text{MB-001, MB-002, MB-003}$ & 3 & 36 & 1.000 & 1.000 & 1.000 & 0.000 & 0.000 \\
\textbf{Cryptographic Key Management} & $\text{MB-004, MB-005, MB-006}$ & 3 & 36 & 1.000 & 1.000 & 1.000 & 0.000 & 0.000 \\
\textbf{Cryptographic Ciphers \& Modes} & $\text{MB-007 \dots MB-012}$ & 6 & 72 & 1.000 & 1.000 & 1.000 & 0.000 & 0.000 \\
\textbf{Local Storage (Preferences)} & $\text{MB-013, MB-014, MB-015}$ & 3 & 36 & 1.000 & 1.000 & 1.000 & 0.000 & 0.000 \\
\textbf{Local Storage (Database \& Files)} & $\text{MB-016, MB-017, MB-018}$ & 3 & 36 & 1.000 & 1.000 & 1.000 & 0.000 & 0.000 \\
\textbf{Manifest \& IPC Configurations} & $\text{MB-019 \dots MB-022}$ & 4 & 48 & 1.000 & 1.000 & 1.000 & 0.000 & 0.000 \\
\textbf{PII Dataflow \& Taint Leaks} & $\text{MB-023 \dots MB-028}$ & 6 & 72 & 1.000 & 1.000 & 1.000 & 0.000 & 0.000 \\
\textbf{Third-Party Telemetry SDKs} & $\text{MB-029, MB-030}$ & 2 & 24 & 1.000 & 1.000 & 1.000 & 0.000 & 0.000 \\
\midrule
\textbf{Global Consolidated Aggregate} & $\mathbf{MB-001 \dots MB-030}$ & \textbf{30} & \textbf{360} & $\mathbf{1.000}$ & $\mathbf{1.000}$ & $\mathbf{1.000}$ & $\mathbf{0.000}$ & $\mathbf{0.000}$ \\
\bottomrule
\end{tabular}%
}
\end{table*}

\subsection{Benchmark Architecture and Catalog}
\label{sec:benchmark:catalog}

The micro-benchmark suite was generated using our automated suite generator (\path{benchmark/apps/generate_benchmark_sources.py}), producing compilable Android projects with isolated manifests, Java classes, and XML configurations:
In network security, $\text{MB-001}$ configures \texttt{network\_security\_config.xml} with pinning; $\text{MB-002}$ permits cleartext HTTP; and $\text{MB-003}$ uses an empty \texttt{checkServerTrusted()} method~\cite{fahl_msc}.
In key management, $\text{MB-004}$ generates AES-256 in \texttt{AndroidKeyStore}; $\text{MB-005}$ uses static hardcoded keys; and $\text{MB-006}$ tests low-entropy constants.
In ciphers and modes, $\text{MB-007}$ uses AES-GCM; $\text{MB-008}$ uses ECB mode; $\text{MB-009}$ uses DES; $\text{MB-010}$ uses RC4; $\text{MB-011}$ uses static IVs; and $\text{MB-012}$ uses predictable PRNG seeds~\cite{crypto_misuse_ccs}.
In storage, $\text{MB-013}$ uses \texttt{EncryptedSharedPreferences}; $\text{MB-014}$ uses plaintext XML; $\text{MB-015}$ uses world-readable storage; $\text{MB-016}$ uses SQLCipher; $\text{MB-017}$ uses plaintext SQLite; and $\text{MB-018}$ writes to external cache.
In manifest/IPC, $\text{MB-019}$ sets \texttt{allowBackup="false"}; $\text{MB-020}$ enables backup; $\text{MB-021}$ enables debugging; and $\text{MB-022}$ exposes exported activities without permissions.
In taint tracking, $\text{MB-023}$ sanitizes PII; $\text{MB-024}$ exfiltrates device IMEI via HTTP; $\text{MB-025}$ persists GPS coordinates; $\text{MB-026}$ logs NID to logcat; $\text{MB-027}$ leaks contacts; and $\text{MB-028}$ intercepts SMS. $\text{MB-029}$ enforces zero trackers, while $\text{MB-030}$ embeds third-party trackers~\cite{binns_third_party_trackers}.

\subsection{Implementation Conformance Verification}
\label{sec:benchmark:results}

We emphasize that Table~\ref{tab:full_benchmark} reports unit-test pass rates on author-constructed exemplars verifying AST visitor implementation correctness, rather than wild-world detection capability (which is evaluated in Section~\ref{sec:ecosystem:adjudication}). We evaluated Reg2App against the complete ground-truth matrix (\path{benchmark/ground_truth.json}), comprising 360 evaluation cells ($30 \text{ apps} \times 12 \text{ rules}$). Across these cells, the ground truth contains 38 active non-conformance defects and 322 conforming controls. As presented in Table~\ref{tab:full_benchmark}, Reg2App correctly detected all 38 non-conformance defects ($TP = 38, FN = 0$) and confirmed all 322 conforming controls ($TN = 322, FP = 0$), yielding defect-detection $\text{Precision} = 1.000$, $\text{Recall} = 1.000$, and $F_1 = 1.000$ with $100\%$ two-class cell agreement ($360/360$). Because the benchmark uses exclusively synthetic apps with no obfuscation, reflection, JNI, or packing, no cells produce $\mathbf{InsufficientEvidence}$; the benchmark verifies that AST visitors correctly implement extraction rules under controlled conditions rather than evaluating evasion resilience. The suite covers 8 of the 12 rules; REV-06 (session binding), REV-08 (anti-tampering), and REV-09 (permission minimization) involve aggregate or behavioral properties not amenable to isolated micro-testcases. Real-world capability on obfuscated production software is evaluated via independent manual adjudication in Section~\ref{sec:ecosystem:adjudication}.

\subsection{Comparative Baseline Analysis: MobSF and Androguard}
\label{sec:benchmark:baseline}

We executed MobSF v3.9.7~\cite{mobsf} and Androguard v4.1~\cite{androguard} across the 30 micro-apps. While both tools were engineered for general vulnerability discovery rather than statutory verification, this comparison highlights the structural mismatch:
(1)~\textit{False Positives}: In $\text{MB-007}$ (secure AES-GCM), MobSF flagged a medium vulnerability because its regex matched \texttt{AES} without parsing the transformation argument ($FPR = 0.167$).
(2)~\textit{False Negatives}: In $\text{MB-014}$, when preference keys were constructed dynamically via string concatenation, MobSF missed the unencrypted write ($FNR = 0.250$).
(3)~\textit{False Compliance Attributions}: In $\text{MB-001}$ (zero vulnerabilities), MobSF emitted an arbitrary ``Security Score: 78/100'', penalizing missing analytics. MobSF conflates absence of features with vulnerability, whereas Reg2App distinguishes mandatory prohibitions from unobservable backend obligations.

\subsection{Component Ablation Study}
\label{sec:benchmark:ablation}

To isolate the empirical contributions of each architectural layer under a common protocol, we evaluated five incremental configurations against the complete 360-cell benchmark suite (Table~\ref{tab:ablation_study}): Config 1 (Ad-Hoc Syntax Baseline), Config 2 (+ 7-Tuple Typed Ontology), Config 3 (+ Observability Operator $\Omega$), Config 4 (+ Evidence Provenance Graph $G$), and Config 5 (Full Reg2App Pipeline). Each configuration was evaluated across six standardized metrics: Precision on detected defects; Recall of ground-truth non-conformance; False Positive Rate on observable code ($FPR_{\text{obs}}$); False Compliance Rate on unobservable obligations ($FCR_{\text{unobs}}$); Provenance Completeness (percentage of verdicts linked to verified AST instruction offsets); and Mean Latency per application.

\begin{table}[b]
\small
\caption{Architectural Ablation Study Across Verification Layers Evaluated on the Common 360-Cell Ground-Truth Benchmark Suite}
\label{tab:ablation_study}
\centering
\resizebox{\columnwidth}{!}{%
\begin{tabular}{lcccccc}
\toprule
\textbf{Configuration Layer} & \textbf{Precision} & \textbf{Recall} & \textbf{$FPR_{\text{obs}}$} & \textbf{$FCR_{\text{unobs}}$} & \textbf{Prov. Comp.} & \textbf{Latency} \\
\midrule
Config 1: Ad-Hoc Syntax Baseline & $66.7\%$ & $73.7\%$ & $16.7\%$ & $33.3\%$ & $0.0\%$ & $12.4 \text{ ms}$ \\
Config 2: + 7-Tuple Typed Ontology & $89.2\%$ & $86.8\%$ & $4.2\%$ & $33.3\%$ & $35.0\%$ & $21.8 \text{ ms}$ \\
Config 3: + Observability Operator $\Omega$ & $94.6\%$ & $92.1\%$ & $1.8\%$ & $\mathbf{0.0\%}$ & $42.0\%$ & $28.5 \text{ ms}$ \\
Config 4: + Evidence Graph $G$ & $\mathbf{100.0\%}$ & $97.4\%$ & $\mathbf{0.0\%}$ & $\mathbf{0.0\%}$ & $\mathbf{100.0\%}$ & $36.1 \text{ ms}$ \\
Config 5: Full Epistemic Pipeline & $\mathbf{100.0\%}$ & $\mathbf{100.0\%}$ & $\mathbf{0.0\%}$ & $\mathbf{0.0\%}$ & $\mathbf{100.0\%}$ & $\mathbf{41.8 \text{ ms}}$ \\
\bottomrule
\end{tabular}%
}
\end{table}

The results demonstrate clear architectural progressions: transitioning from unstructured regex matching to the typed 7-tuple ontology (Config 2) cuts $FPR_{\text{obs}}$ from $16.7\%$ to $4.2\%$. Introducing the Observability Operator $\Omega$ (Config 3) deterministically eliminates false compliance attributions on unobservable backend/organizational obligations ($FCR_{\text{unobs}}: 33.3\% \rightarrow 0.0\%$), an invariant guaranteed by construction via Invariant~1. Constructing the Evidence Provenance Graph $G$ (Config 4) eliminates remaining false alarms by demanding affirmative witness subgraphs, achieving $100\%$ provenance completeness. Finally, incorporating the full epistemic assessment space (Config 5) integrates bounded semantics over partial observations, reaching $100\%$ recall across all benchmark specifications. Developer remediation efficiency across human reporting formats is independently evaluated via the Latin Square experiment in Section~\ref{sec:usable}.

\subsection{Computational Overhead}
\label{sec:benchmark:overhead}

Benchmarked on an AMD Ryzen 7 workstation (32 GB RAM), mean analysis time was $41.8 \pm 4.2 \text{ ms}$ per micro-app ($1.25 \text{ s}$ total for 360 cells), with memory usage under $85 \text{ MB}$, confirming feasibility for CI/CD integration.
